% TOQB note: calibration-only placement against self-measuring placement (workplace, 2 November 2026).
% Not declared; not a paper. Whether it goes into Paper 6 as an appendix is decided later by home and the owner.
% Every simulation number carries "% src: out/construction_result.json <key>" (written by construction_sim.py).
\documentclass[10pt,a4paper]{article}
\usepackage[margin=2.2cm]{geometry}
\usepackage{amsmath,amssymb,amsthm}
\usepackage{booktabs}
\usepackage[hidelinks]{hyperref}
\usepackage{url}
\setlength{\parskip}{4pt}
\setlength{\parindent}{0pt}
\sloppy
\newtheorem{lemma}{Lemma}
\newtheorem{proposition}{Proposition}
\newtheorem{corollary}{Corollary}
\theoremstyle{definition}
\newtheorem{definition}{Definition}
\theoremstyle{remark}
\newtheorem*{remark}{Remark}

\title{Calibration-only placement against self-measuring placement:\\ a near-trivial lower bound, a construction,
and what neither proves\\ \large TOQB note, draft (not declared)}
\author{The Architect}
\date{2 November 2026 (CET)}

\begin{document}
\maketitle

\begin{abstract}
A placement policy that reads nothing about the device except its coupling graph and its calibration numbers gives
the same placement law on any two device states with the same calibration numbers. This is close to a tautology, and
we state it only to fix terms (Lemma~1). The content of the note is elsewhere: (i) an explicit pair of noise models
on IBM's FakeKingston coupling graph, an error that acts only when couplers at distance one are driven in the same
layer, under which every isolated two-qubit benchmark returns identical counts while a ring echo through the planted
coupler dies (Section~4, checked in Aer); and (ii) a list of what the lemma and the matching upper bound for a
self-measuring policy do not prove (Section~5). Nothing here is a device result.
\end{abstract}

\section{Terms}

\begin{description}
\item[$\theta$] the device's true state: every element's noise, including errors that appear only when several
  elements are driven together, at the time of the run.
\item[$C(\theta)$] the calibration numbers. We take them to be static per-element numbers (coherence times, readout)
  together with the outcome distributions of a fixed set $\mathcal E_{\rm iso}$ of \emph{isolated} experiments:
  experiments in which no two couplers at distance one carry a two-qubit gate in the same layer. A randomized
  benchmark of one coupler is such an experiment; a layer-fidelity measurement on a chain of simultaneous gates is
  not.
\item[policy $\pi$] a map from (working coupling graph $G$, $C(\theta)$) to a probability law over placements $L$ of a
  ring of $n$ qubits. $\Pi_C$ is the class of such maps. It contains every compiler pass that decides from those inputs
  alone, deterministic or randomized, open or closed source. It does \emph{not} contain policies that read numbers of
  simultaneously driven layers, policies that read results of the user's own circuits, or policies that run their own
  test circuits on the day (the self-measuring class of Section~3).
\item[regret $R_\theta(L)$] how far the result of placement $L$ falls short of a comparator placement in state
  $\theta$. The comparator is fixed in (A3) below; it is not ``the best placement on the device''.
\end{description}

\section{The lower bound (two-point / averaging)}

\begin{lemma}\label{lem:two}
Let $\theta_1,\dots,\theta_M$ be states with $C(\theta_k)=c$ for all $k$, and $\pi\in\Pi_C$. Then
\[
\max_k \mathbb E_{L\sim\pi(G,c)}\, R_{\theta_k}(L)\;\ge\;\frac1M\,\min_L \sum_{k=1}^M R_{\theta_k}(L).
\]
\end{lemma}
\begin{proof}
The law $\mu=\pi(G,c)$ is the same for every $k$, so
$\max_k \mathbb E_\mu R_{\theta_k} \ge \frac1M\sum_k \mathbb E_\mu R_{\theta_k} = \mathbb E_\mu\big[\frac1M\sum_k
R_{\theta_k}(L)\big] \ge \frac1M\min_L\sum_k R_{\theta_k}(L)$.
\end{proof}

This is the two-point method of Le Cam in its simplest averaged form (see Yu~\cite{yu1997} for the method and its
relatives). The lemma says nothing until a family $\{\theta_k\}$ with equal $C$ and a positive right-hand side is
exhibited; with $M=2$ and one placement that is good in both states the right-hand side is $0$. Section~4 gives such
a family.

\begin{corollary}\label{cor:ring}
Let $E_w$ be the working couplers whose two qubits work, and for $e\in E_w$ let $\theta_e$ be the baseline state
plus the same-layer error of Section~4 on $e$. Assume (A2) and (A3) of Section~3, with a cost $\Delta>0$ for a ring through $e$. Every
ring placement of $n$ qubits uses $n$ distinct couplers of $E_w$, so
$\frac{1}{|E_w|}\min_L\sum_e R_{\theta_e}(L)\ge n\Delta/|E_w|$, and every $\pi\in\Pi_C$ has
$\max_e \mathbb E\,R_{\theta_e}\ge n\Delta/|E_w|$. On FakeKingston $|E_w|=164$: the bound is $0.12\,\Delta$,
$0.24\,\Delta$ and $0.44\,\Delta$ for $n=20$, 40 and 72.
\end{corollary}
% src: |E_w| = 164 counted with psf_os.device (working couplers with both qubits ok, FakeKingston); 20/164, 40/164, 72/164
The step $C(\theta_e)=C(\theta_0)$ is exact, not approximate: the planted error is inserted only after a CZ on $e$
whose layer also drives a coupler at distance one from $e$, and an isolated experiment contains no such layer, so its
circuit and hence its outcome law are unchanged (V1 below confirms this numerically).

\section{The upper bound for a self-measuring policy, with its assumptions}

The self-measuring policy is psf-os's chain: a scan (the mirror echo at $D=24$, read per ring site), a registry
(a path is listed after two jobs agree), and a placement by the same fixed rule with the middle coupler of each listed
path masked. The bound below holds only under these assumptions:
\begin{itemize}
\item[(A1)] Where the echo is dead, the user's circuit is bad too. (Paper 6's forward test on \texttt{ibm\_marrakesh}
  is consistent with this: the gain sat on the replaced qubits. It is not a proof.)
\item[(A2)] The result of a placement depends only on whether its ring contains a dead path: cost $\Delta$ if it
  does, $0$ if not. This is a simplification; real rings differ in many other ways.
\item[(A3)] The comparator is not the best placement on the device. It is the placement the rule makes when every
  truly dead path is masked, and we assume the rule can place a ring with those masks (psf-os does not search for the
  best placement).
\item[(A4)] The state does not change between the scan and the user's job. A gray (weak, same-day) entry breaks this
  by definition; the sentinel (M6) exists to detect such a change inside the user's job.
\end{itemize}

\begin{definition}
With the dead line $\tau=0.12$, a site is \emph{$t$-dead} if its true signal $a_j\le\tau-t$ and \emph{$t$-healthy} if
$a_j\ge\tau+t$. A ring is \emph{$t$-separated} if every site is one or the other.
\end{definition}

\begin{proposition}\label{prop:upper}
A site signal is $\hat a_j=\frac14\sum_{i=1}^4 \bar m_{i,j}$, each $\bar m_{i,j}$ the mean of $N$ independent outcomes
in $\{-1,+1\}$ (four inputs, $N$ shots each). By Hoeffding's inequality~\cite{hoeffding1963} for the $4N$ bounded
outcomes, $\Pr(|\hat a_j-a_j|\ge t)\le 2e^{-2Nt^2}$. On a $t$-separated ring of $n$ sites, one scan puts every site on
the correct side of $\tau$ with probability at least $1-2n\,e^{-2Nt^2}$, and then its dead runs are exactly the true
ones. Under (A1)--(A4) the expected regret of the policy is at most $4n\Delta\,e^{-2Nt^2}$ (a true dead path is missed
only if one of the two scans misreads; a healthy path is listed only if both do, which under (A2) costs nothing as long
as (A3) holds with the extra mask). Regret $\le\varepsilon$ needs $N\ge \ln(4n\Delta/\varepsilon)/(2t^2)$ shots per
input, logarithmic in $n$.
\end{proposition}
% src: Hoeffding with 4N samples of range 2: exp(-2 (4N)^2 t^2 / (4N * 4)) = exp(-2 N t^2)

\textbf{Numbers.} With $N=1000$: a site with true signal $0$ is read above $\tau$ with probability at most
$e^{-28.8}\approx3\times10^{-13}$ (two-sided $6\times10^{-13}$); for $n=72$ the regret bound is
$\approx9\times10^{-11}\Delta$. A site at the device's healthy level (0.35--0.45, so $t\ge0.23$) is read dead with
probability below $10^{-45}$. The earlier estimate of $0.08\,\%$ counted one input of $1000$ shots ($e^{-7.2}$); it is
conservative by a factor of four in the exponent.
% src: exp(-2*1000*0.12^2) = exp(-28.8) = 3.1e-13; 4*72*exp(-28.8) = 8.9e-11; 2*exp(-2*1000*0.23^2) = 2.3e-46; exp(-7.2) = 7.5e-4

\textbf{What these numbers are not.} The line $\tau=0.12$, the three sites and the two jobs are choices made for
psf-os, not outputs of the statistics. The bound says how often a reading lands on the wrong side of a chosen line,
given a margin $t$. A site whose true signal sits near the line has no useful margin: at $a_j=0.10$, $t=0.02$ and the
bound is $2e^{-0.8}\approx0.90$. Section~4 meets exactly this case.

\section{The construction (the centre of the note)}

\textbf{State A.} On FakeKingston, after each CZ on the planted coupler $e$, if a coupler at distance one from $e$
(disjoint from $e$, joined to it by one coupler) carries a CZ in the same ring layer, an extra ZZ error acts on $e$.
The ring layer of a CZ is read from the compiled circuit by occurrence (the $k$-th CZ on a coupler is in layer $k$;
every ring coupler carries one CZ per layer, asserted). Two models, fixed before any run and not tuned: coherent
$R_{ZZ}(\pi/8)$, run as a statevector on $n=20$ rings with the scan's own echo ($h=0.6$); and Pauli $ZZ$ with
probability $0.15$ (the value of psf-os test T3), run on the stabilizer method with the echo at its Clifford point
($h=\pi/4$, as T3), so that $n=40$ and 72 can be simulated. \textbf{State B} is the same device with no planted error.
No other noise is modelled, so B reads $1.000$ everywhere; that is not a realistic level, and only differences between
A and B carry meaning.

$e=(42,43)$ is fixed by rule: positions 12 and 13 (the positions at which psf-os test T3 plants on the $n=40$ ring), taken here on the standard $n=20$ ring. It lies on the
standard rings of $n=20$, 40 and 72. On \texttt{ibm\_kingston} these qubits were healthy (Paper 6); the choice does not
imitate the device's dead block. Its couplers at distance one are $(36,41)$, $(40,41)$, $(44,45)$, $(56,63)$.
% src: out/construction_result.json e, distance_one

\begin{table}[h]\centering\small
\begin{tabular}{p{0.36\textwidth}p{0.42\textwidth}p{0.14\textwidth}}
\toprule
check & state A & state B\\
\midrule
V1 benchmark of $e$ alone, $m=2,8,32$ (both models) & 0 errors inserted; counts identical to B & success 1.000\\
V1 the same with $(36,41)$ in the same layer, Pauli & success 0.74625, 0.2685, 0.24975 & 1.000\\
V1 the same, coherent & success 0.8755, 0.1745, 0.0250 & 1.000\\
\midrule
V2 Pauli scan, $n=20/40/72$: dead run over $e$ & 16 sites each; usable 4, 24, 56 & usable 20, 40, 72\\
V2 coherent scan, $n=20$ & dead run of 3 sites (44, 43, 42), mean 0.10; ring mean 0.39 & 1.000\\
\midrule
V3 Pauli: listed path; usable, standard $\to$ placed ring & $[43,42,41]$; $24\to40$ ($n=40$), $56\to72$ ($n=72$) & --\\
\quad mean signal on sites already usable & $0.83\to1.00$, $0.93\to1.00$; sentinel CLEAN & --\\
V3 coherent ($n=20$): listed path; usable & $[43,42]$; $20\to20$ & --\\
\quad mean signal over the ring & $0.40\to1.00$; sentinel CLEAN & --\\
\midrule
V4 coherent, $n=20$, no flip & usable 16, dead run (44, 43, 42, 41); mean 0.40 & 1.000\\
V4 flip at position 12 / at 13 / global (even positions) & usable 20, every site 1.000 & 1.000\\
\bottomrule
\end{tabular}
% src: out/construction_result.json V1 (success_A, success_B, inserted, identical_counts); V2_pauli.*.dead/usable;
%      V2_coherent.dead (mean 0.1025), mean(a_scan1) = 0.392; V3_pauli.listed, user.40/72 usable, mean_usable;
%      V3_coherent.listed, user standard/avoiding usable 20/20, mean_all 0.402/1.000; V4 rows
\caption{Aer, FakeKingston, planted same-layer error on $e=(42,43)$. Simulation, not a device result.}
\label{tab:sim}
\end{table}

\textbf{Reading the table.} V1 is the step the lemma needs: any number made from experiments on $e$ alone is the same
in A and B, exactly (no gate is inserted), while a benchmark that drives a distance-one coupler in the same layer sees
the error at once. So the construction stands only against calibration that is isolated in our sense; against a
layer-fidelity measurement that covers $e$ and a neighbour it fails, which is why such policies are outside $\Pi_C$.
V2 and V3 show the two sides of the separation: a policy in $\Pi_C$ places the same standard ring in A and B, and in A
that ring is dead around $e$; psf-os's scan lists the path after two scans and the placed ring reads clean.

The headline and the essential number differ in the coherent case. The headline is ``usable sites $20\to20$'', no
gain in width. The essential number is the ring mean, $0.40\to1.00$: the coherent error lowered the whole ring, but
only three sites fell under the line, at a mean of 0.10, inside the margin where Proposition~\ref{prop:upper} says
nothing. The same circuit was flagged dead in three of four runs (both scans, V4's unflipped run) and not in one (the
standard user run). That psf-os listed the path here was, in that sense, luck of the shots.
% src: V2_coherent scan 1 dead (11,3); scan 2 dead (12,3) (re-read from the work folder); V3_coherent standard dead [] ; V4 no flip dead [[44,43,42,41]]
In the Pauli model the dead run is 16 sites wide at every size: at the Clifford point a Pauli error spreads at full
speed along the ring, so the width depends on the model as much as on the error.

V4: the coherent error is a static ZZ on one coupler, and a local frame flip at either end of $e$ (or the global flip)
cancels it exactly in this mirror circuit, as psf-os's repair notes state. The model has no other error, so the flip
costs nothing here; on the device the global flip cost the healthy part of the ring.

\section{What neither the lemma nor the bound proves}

\begin{itemize}
\item \textbf{Whether such states occur, and how often.} The lemma needs states with equal $C$ and different damage;
  the construction shows they can exist. How often a real device is in one is a question of observation. Our evidence
  is \texttt{ibm\_kingston} on three days, \texttt{ibm\_fez} and \texttt{ibm\_marrakesh} on one day each.
\item \textbf{(A1), the echo against the user's circuit.} Assumed, not shown.
\item \textbf{(A2) and (A3).} The cost model is binary and the comparator is the rule's own family. Neither says
  psf-os's placement is the best one on the device.
\item \textbf{(A4).} A change between scan and job breaks the bound; the sentinel is meant to catch it, and it has
  only run in simulation.
\item \textbf{Scope.} ``Any company's compiler'' is too broad. The lemma covers policies that read only the coupling
  graph and isolated calibration numbers. A policy that reads simultaneous-layer numbers (for example layer fidelity on
  a chain) or the user's own results is outside it, and the construction must then move to couplers those numbers do
  not cover, if any.
\item \textbf{The thresholds.} $\tau=0.12$, three sites and two jobs are chosen values. Near the line the reading is a
  coin toss (Section~4, coherent case).
\item \textbf{The model's width.} The size of a dead run in simulation depends on the model (3 sites coherent, 16
  sites Pauli at the Clifford point). Neither is a claim about the device.
\end{itemize}

\section{Relation to the device observation}

On \texttt{ibm\_kingston} (30 October to 1 November 2026) the same 18 qubits of the standard $n=40$ ring formed a dead
run in every run of a byte-identical circuit; the calibration showed nothing there~\cite{paper5}; the block behaved in
part coherently (twirling turned its inversion into noise; its signal crossed zero over depth); and a frame flip on
qubits 27 and 29 lifted it from $-0.19$ to $+0.13$ (31 October) and from $-0.19$ to $+0.11$ (1 November) at $D=24$ (Paper 6, draft). The construction is
\emph{consistent} with these features: an error invisible to isolated calibration, coherent, and cancelled in a mirror
circuit by a local flip. It is not shown to be their cause, and two features differ: on the device the flip repaired
the block only in part, and the global flip cost the healthy part of the ring; in the model both flips repair fully
and cost nothing.

\begin{thebibliography}{9}
\bibitem{yu1997} B.~Yu, ``Assouad, Fano, and Le Cam,'' in \emph{Festschrift for Lucien Le Cam}, Springer, New York,
  1997. doi:10.1007/978-1-4612-1880-7\_29. [Page range: CHECK.]
\bibitem{hoeffding1963} W.~Hoeffding, ``Probability inequalities for sums of bounded random variables,''
  \emph{J. Amer. Statist. Assoc.} \textbf{58}(301), 13--30 (1963).
\bibitem{paper5} The Architect, ``Readable Depth Without the Calibration Map'' (Paper 5),
  doi:10.5281/zenodo.23827536.
\end{thebibliography}

\end{document}
